\documentclass[10pt,a4paper]{amsart}
\usepackage[margin=19mm]{geometry}
\usepackage{amsmath,amssymb,mathtools}
\usepackage[scr=rsfs,cal=euler]{mathalfa}
\usepackage{xfrac}
\usepackage[unicode,hidelinks,bookmarks=false]{hyperref}
\newcommand{\ie}{i.e.\ }
\newcommand{\cf}{cf.\ }
\newcommand{\resp}{respectively\ }
\newcommand{\Subsection}{\subsection}
\providecommand{\nobreakdash}{\nobreak\hbox{-}}
\setlength{\emergencystretch}{2em}
\multlinegap0pt
\usepackage{amssymb}
\newtheorem{lemma}{Lemma}
\title{Sharp inf-sup estimate for the Stokes equation in tight domains with periodic pillars and some numerical implications}
\author{Qi Xin, Shihua Gong, Jinchao Xu}
\date{Source: arXiv:2604.12643v2 (CC BY 4.0)}
\begin{document}
\maketitle

\noindent\emph{Source and licence.} Sharp inf-sup estimate for the Stokes equation in tight domains with periodic pillars and some numerical implications. Version arXiv:2604.12643v2 (CC BY 4.0), distributed by arXiv under the Creative Commons Attribution 4.0 International licence, http://creativecommons.org/licenses/by/4.0/, as recorded in the arXiv licence metadata for this exact version and shown on the abstract page https://arxiv.org/abs/2604.12643v2. Full licence notice: \texttt{licenses/arxiv-2604.12643-CC-BY-4.0.txt}.\par\smallskip

\noindent\textit{Editorial scope.} Counted unit: the article's lemma lemma:poincare (source0.tex lines 232--237). The article's complete original proof of it is delivered with it (source0.tex lines 239--267, one \texttt{proof} environment of 29 lines). No mathematics is written, repaired or re-derived: the statement and the proof are the article's own lines, in the article's own notation, and no retained line is substituted or re-flowed.  Retained uncounted as the source-local context the proof needs: lines 144--144.  Nothing was omitted from any retained span, so the package carries no omission record.  No retained line contains a citation, so the wrapper carries no bibliography and no bibliography entry.  No retained line contains a figure environment, an image-inclusion command or drawing code, so no image is delivered and the package declares no asset dependency.  Editorial formatting only, in addition: an AMS \texttt{amsart} preamble that restates the article's own declarations and macros used by the retained lines, namely the environments \texttt{lemma} on separate counters, together with the standard packages those lines need.  The retained environments print in this document as Lemma 1 in order of appearance, whereas the article prints the same items with the numbering of its own full document.  The complete native source of the article as distributed in the arXiv e-print for this exact version is bundled unchanged as \texttt{sources/198414/source0.tex}.
\section{Mathematical Model and Main Results}\label{sec:problem_main_results}

\begin{lemma}[Global Poincaré Inequality]\label{lemma:poincare}
    Let $\Omega_m$ be the fluid domain with dense obstacles defined in Section \ref{sec:problem_main_results}, with zero boundary conditions on the holes $\partial T_{\epsilon}^{i,j}$ and the exterior boundary $\partial \Omega$. There exists a constant $C > 0$, independent of $\epsilon$ (and thus independent of the obstacle density), such that for all $\mathbf{v} \in H_0^1(\Omega_m)^d$:
    \begin{equation}
        \|\mathbf{v}\|_{L^2(\Omega_m)} \le C \min\{1, \sigma_{\epsilon}\} \|\nabla \mathbf{v}\|_{L^2(\Omega_m)}.
    \end{equation}
\end{lemma}

\begin{proof}
    We extend $\mathbf{v} \in H_0^1(\Omega_m)^d$ by zero inside the obstacles $T_{\epsilon}^{i,j}$ to obtain $\tilde{\mathbf{v}} \in H_0^1(\Omega)^d$. The classical Poincaré inequality on the macroscopic domain $\Omega$ immediately yields $\|\tilde{\mathbf{v}}\|_{L^2(\Omega)} \le C_P \|\nabla \tilde{\mathbf{v}}\|_{L^2(\Omega)}$, providing the global upper bound $C_P$ independent of $\epsilon$.

    To derive the $\sigma_\epsilon$-dependent local bound, we decompose $\Omega$ into the union of active fluid cells $P_\epsilon^{i,j}$ and a boundary layer $\Lambda_{\epsilon}$ near $\partial \Omega$ containing no complete obstacles. The width of $\Lambda_{\epsilon}$ is of order $\mathcal{O}(\epsilon)$. Since $\mathbf{v} = 0$ on $\partial \Omega$, the standard one-dimensional Poincaré inequality across the thickness of $\Lambda_{\epsilon}$ gives
    \begin{equation}
        \|\mathbf{v}\|_{L^2(\Lambda_{\epsilon})}^2 \le C_0 \epsilon^2 \|\nabla \mathbf{v}\|_{L^2(\Lambda_{\epsilon})}^2.
    \end{equation}

    For each cell $P_\epsilon^{i,j}$, let $B_{\mathrm{out}}^{i,j}$ be its circumscribed ball with radius $R_{\epsilon} = C_{\mathrm{out}} \epsilon$, and $B_{\mathrm{in}}^{i,j} \subset T_{\epsilon}^{i,j}$ be the inscribed ball with radius $r_{\epsilon} = \alpha a_{\epsilon}$ (where $\alpha, C_{\mathrm{out}} > 0$ are geometric lattice constants). Since $\mathbf{v} = 0$ on $\partial T_{\epsilon}^{i,j}$, it vanishes on $\partial B_{\mathrm{in}}^{i,j}$. Using polar coordinates $(r, \theta)$ centered at the obstacle, the fundamental theorem of calculus along radial segments yields
    \begin{equation}
        \mathbf{v}(r, \theta) = \int_{r_{\epsilon}}^r \frac{\partial \mathbf{v}}{\partial \rho}(\rho, \theta) \, \mathrm{d}\rho.
    \end{equation}
    Applying the Cauchy-Schwarz inequality, we separate the derivative and the geometric weight:
    \begin{equation}
        |\mathbf{v}(r, \theta)|^2 \le \left( \int_{r_{\epsilon}}^r \left| \frac{\partial \mathbf{v}}{\partial \rho} \right|^2 \rho \, \mathrm{d}\rho \right) \left( \int_{r_{\epsilon}}^r \frac{1}{\rho} \, \mathrm{d}\rho \right) \le \ln\left(\frac{r}{r_{\epsilon}}\right) \int_{r_{\epsilon}}^{R_{\epsilon}} |\nabla \mathbf{v}|^2 \rho \, \mathrm{d}\rho.
    \end{equation}
    Integrating this point-wise bound over the outer ball $B_{\mathrm{out}}^{i,j}$, we obtain
    \begin{align}
        \|\mathbf{v}\|_{L^2(B_{\mathrm{out}}^{i,j})}^2 = \int_0^{2\pi} \int_{r_{\epsilon}}^{R_{\epsilon}} |\mathbf{v}(r, \theta)|^2 r \, \mathrm{d}r \mathrm{d}\theta & \le \|\nabla \mathbf{v}\|_{L^2(B_{\mathrm{out}}^{i,j})}^2 \int_{r_{\epsilon}}^{R_{\epsilon}} r \ln\left(\frac{r}{r_{\epsilon}}\right) \, \mathrm{d}r \nonumber                                 \\
                                                                                                                                                                      & = \|\nabla \mathbf{v}\|_{L^2(B_{\mathrm{out}}^{i,j})}^2 \left[ \frac{R_{\epsilon}^2}{2} \ln\left(\frac{R_{\epsilon}}{r_{\epsilon}}\right) - \frac{R_{\epsilon}^2 - r_{\epsilon}^2}{4} \right].
    \end{align}
    Dropping the negative term and substituting the radii $R_{\epsilon} = C_{\mathrm{out}}\epsilon$ and $r_{\epsilon} = \alpha a_{\epsilon}$, the geometric bracket is bounded by $\frac{1}{2} C_{\mathrm{out}}^2 \epsilon^2 \ln\left(\frac{C_{\mathrm{out}} \epsilon}{\alpha a_{\epsilon}}\right)$. By defining the characteristic scaling $\sigma_{\epsilon}^2 := \epsilon^2 |\ln(a_{\epsilon}/\epsilon)|$, there exists a constant $C_1 > 0$ such that the integral evaluates to $\le C_1 \sigma_{\epsilon}^2$.

    Summing over all cells $P_\epsilon^{i,j}$ and accounting for the finite overlap of the circumscribed balls, we have $\sum_{i,j} \|\mathbf{v}\|_{L^2(P_\epsilon^{i,j})}^2 \le C_2 \sigma_{\epsilon}^2 \|\nabla \mathbf{v}\|_{L^2(\Omega_m)}^2$. Combining this with the boundary layer bound, we conclude
    \begin{equation}
        \|\mathbf{v}\|_{L^2(\Omega_m)}^2 = \|\mathbf{v}\|_{L^2(\Lambda_{\epsilon})}^2 + \sum_{i,j} \|\mathbf{v}\|_{L^2(P_\epsilon^{i,j})}^2 \le C_3 \sigma_{\epsilon}^2 \|\nabla \mathbf{v}\|_{L^2(\Omega_m)}^2.
    \end{equation}
    Taking the square root and combining this microscopic limit with the macroscopic bound $\|\mathbf{v}\|_{L^2(\Omega_m)} \le C_P \|\nabla \mathbf{v}\|_{L^2(\Omega_m)}$ completes the proof.
\end{proof}

\end{document}
